\chapter{Methodik / Entwicklung der Methodik}
\label{chap:Principles}
Der Hauptteil der Arbeit gliedert sich in zwei Teile. Im ersten Teil (Methodik bzw. Entwicklung der Methodik) wird die grundlegende Methodik der Arbeit erarbeitet. Dabei werden grundlegende Fragen erläutert (Wie gehe ich vor? Was ist bekannt? Welche Annahmen werden getroffen? etc.), wobei die übergeordnete Forschungshypothese (bei einer MA) bzw. Zielsetzung der Arbeit (bei einer BA, PA, MT) berücksichtigt wird. Eine Erläuterung der Grundlagen ist nur da erlaubt, wo sie die täglichen Arbeiten eines Ingenieury der Strukturmechanik übersteigt. Der Umfang des Fachwissen kann sich in der Auswahl der Methodik, Interpretation der Ergebnisse und der kritischen Reflexion der gewählten Vorgehensweise widerfinden. Dieses Kapitel stellt die Grundlage der gesamten Arbeit dar.\\
Hier gelten zwei Regeln:
\begin{itemize}
	\item Von Experten - Für Experten!
	\item So wenig wie möglich und so viel wie nötig!
\end{itemize}

Blindtext: Nfeste his Questus mox se opportunitatus sto appropinquo alica distinguo nutus tutela pio Suffusus si hic exesto tristis Seorsum, to diu Nitor qua Irrisorie ora Orexis. Tutus infervesco Editio saeta his Luctus, his apud Grator manus Edico hic Exsupero libens tumultuarius, bos satago edo to Hinc diligentia Inflo lea ago hac mores Vergo dux Renovatio letalis. No Declino vir excito utercumque Percut.
\section{Erstes Unterkapitel der Theorie}
\label{sec:Theorie1}
Blindtext: Acto re stupeo Labor sus, ver ex aut exhorto sis aliter foetidus expono. Sensus apud latrocinor, impenetrabiilis far incrementabiliter Commodo cum mel voluptarius Pariter modicus opto coepto, maligo spes Resono Curvo escendo adsum per Frutex, ubi ait animadverto poema, adicio Consonum archipater sum Aeger Dux prius edo paterna precipue, cunae declaratio per dolositas Huic quod Sis canalis quam nam fio Insidiae, si pax Cupido, ut Tergo, ac Cui per quo processus Disputo sui Infucatus leo, ait ops, duo Prodoceo par Verber, nec Uberrime alo Scelestus, res Tellus mei Escensio Mundus, ita liber qui has inconsideratus nauta effrenus, Algor infrunitus, inconcussus Rogo eo non Namucense, commissum, laureatus Scutum, de boo si anhelo Commoneo procellosus sono emitto Crimen agna. Si subo Accubo castimonia hic ibi qua lux sto eu Pulcher Sem. Dis Cubiculum quo scitus Litigo diripio ango quies pes res penitentia Tabula, vos diu Sordes vae Epulor ile Tenor, nox Opulentia diu, ago Suppono sto pia Eri. \gls{FCS}